% GENERATED FILE. Edit canonical lessons, then run npm run generate:books.
% canonical-source-hash: fnv1a64:2a1545add122322e
% canonical-lessons: KA-C315-khusi, KA-C315-santosavayitu, KA-C315-guri, KA-C315-hogabeku-anta-ase, KA-R315-feelings-and-hopes

\chapter{Feelings and Hopes}
\label{ch:ka-feelings-and-hopes}
\hlchaptermodality{315}

\begin{chapteropening}
\textbf{By the end of this chapter:} \emph{I can say how I felt with nanage … āyitu, how I am with -āgiddēne, name a goal with guri, and say what I would love to do or become with -bēku anta āse ide.}

Complete the last lesson of chapter 315: say how you felt and how you are, and what you hope to do or become.
\end{chapteropening}

\section[\texorpdfstring{\kn{ಖುಷಿ} (khuṣi)}{ಖುಷಿ (khuṣi)}]{\kn{ಖುಷಿ} (khuṣi) — gladness, happiness}

\label{lesson:KA-C315-khusi}



\begin{quote}
Before the new one: say \textquotedblleft{}joy\textquotedblright{}, \textquotedblleft{}sadness\textquotedblright{} and \textquotedblleft{}fear\textquotedblright{}. (\emph{santōṣa}, \emph{duḥkha}, \emph{bhaya}.) Then \textquotedblleft{}My bag is lost.\textquotedblright{} (\emph{nanna cīla kaḷeduhōyitu}.)
\end{quote}

\subsection*{You'll want to know: \kn{ಖುಷಿ}}
\textbf{\kn{ಖುಷಿ}} — \emph{khuṣi} — \textquotedblleft{}gladness, happiness\textquotedblright{}.

\begin{quote}
\textbf{\kn{ನಿಮ್ಮನ್ನು} \kn{ನೋಡಿ} \kn{ನನಗೆ} \kn{ತುಂಬಾ} \kn{ಖುಷಿ} \kn{ಆಯಿತು}.}
\end{quote}

\begin{quote}
\emph{nimmannu nōḍi nanage tumbā khuṣi āyitu.} — \textquotedblleft{}I was very glad to see you.\textquotedblright{}
\end{quote}

\textbf{\kn{ನಿಮ್ಮನ್ನು}} (\emph{nimmannu}, \textquotedblleft{}you\textquotedblright{}), as the one seen; \textbf{\kn{ನೋಡಿ}}, \textquotedblleft{}having seen\textquotedblright{}. The gladness \textquotedblleft{}happened\textquotedblright{} to me: \textbf{\kn{ನನಗೆ} … \kn{ಖುಷಿ} \kn{ಆಯಿತು}}, with the \textbf{\kn{ನನಗೆ}} of \textbf{\kn{ನನಗೆ} \kn{ಕನ್ನಡ} \kn{ಗೊತ್ತು}}.

\textbf{\kn{ಖುಷಿ}} and \textbf{\kn{ಸಂತೋಷ}} both mean happiness. \textbf{\kn{ಸಂತೋಷ}} is also the word you say on meeting someone; in everyday talk about how you feel you will hear \textbf{\kn{ಖುಷಿ}} a great deal.

\begin{cousinweb}
\textbf{\kn{ಖುಷಿ}} (\emph{khuṣi}) is a loan, usually traced to Persian \emph{khushī}, \textquotedblleft{}happiness\textquotedblright{}, from \emph{khush}, \textquotedblleft{}glad, pleasant\textquotedblright{}. The same word is everyday Hindi and Urdu. It is not the only Persian piece in words you know: the \emph{havā} of \textbf{\kn{ಹವಾಮಾನ}}, \textquotedblleft{}weather\textquotedblright{}, is Persian for \textquotedblleft{}air\textquotedblright{}.
\end{cousinweb}

\subsection*{Your turn}
Say these aloud:

\begin{itemize}
\raggedright
  \item \emph{khuṣi}, slowly, then at speed
  \item \textquotedblleft{}I was very glad to see you\textquotedblright{} — \emph{nimmannu nōḍi nanage tumbā khuṣi āyitu}
  \item \textquotedblleft{}I was glad\textquotedblright{} — \emph{nanage khuṣi āyitu}
\end{itemize}

\begin{tcolorbox}[breakable,colback=teal!4,colframe=teal!35!black,title={Before you move on}]
What does \textbf{\kn{ಖುಷಿ}} mean? (\textbf{Gladness, happiness.}) \textquotedblleft{}I was glad\textquotedblright{}? (\textbf{\kn{ನನಗೆ} \kn{ಖುಷಿ} \kn{ಆಯಿತು}}.) Who takes \textbf{\kn{ನನಗೆ}}? (\textbf{The one the feeling happens to.})
\end{tcolorbox}
\section[\texorpdfstring{\kn{ನನಗೆ} \kn{ಸಂತೋಷವಾಯಿತು} (nanage santōṣavāyitu)}{ನನಗೆ ಸಂತೋಷವಾಯಿತು (nanage santōṣavāyitu)}]{\kn{ನನಗೆ} \kn{ಸಂತೋಷವಾಯಿತು} (nanage santōṣavāyitu) — how you felt}

\label{lesson:KA-C315-santosavayitu}



\begin{quote}
Say \textquotedblleft{}I was very glad to see you\textquotedblright{}. (\emph{nimmannu nōḍi nanage tumbā khuṣi āyitu}.) Then \textquotedblleft{}I know Kannada\textquotedblright{}. (\emph{nanage kannaḍa gottu}.)
\end{quote}

\subsection*{You'll want to know: a feeling that comes to you}
\begin{quote}
\textbf{\kn{ನಿನ್ನೆ} \kn{ನನಗೆ} \kn{ತುಂಬಾ} \kn{ಸಂತೋಷವಾಯಿತು}.}
\end{quote}

\begin{quote}
\emph{ninne nanage tumbā santōṣavāyitu.} — \textquotedblleft{}Yesterday I was very happy.\textquotedblright{}
\end{quote}

Kannada builds most feelings the way it builds \textbf{\kn{ನನಗೆ} \kn{ಕನ್ನಡ} \kn{ಗೊತ್ತು}}: the person takes \textbf{\kn{ನನಗೆ}}, \textquotedblleft{}to me\textquotedblright{}, and the feeling happens to them. \textbf{\kn{ಸಂತೋಷ}} + \textbf{\kn{ಆಯಿತು}} $\to$ \textbf{\kn{ಸಂತೋಷವಾಯಿತು}}, with the \textbf{-\kn{ವ}-} you met in \textbf{\kn{ಸ್ವಚ್ಛವಾಗಿದೆ}}. Any feeling word fits the frame:

\begin{itemize}
\raggedright
  \item \textbf{\kn{ನನಗೆ} \kn{ಭಯವಾಯಿತು}} (\emph{nanage bhayavāyitu}) — I got scared
  \item \textbf{\kn{ನನಗೆ} \kn{ಬೇಸರವಾಯಿತು}} (\emph{nanage bēsaravāyitu}) — I felt low; I got bored
  \item \textbf{\kn{ನನಗೆ} \kn{ದುಃಖವಾಯಿತು}} (\emph{nanage duḥkhavāyitu}) — I felt sad
  \item \textbf{\kn{ನನಗೆ} \kn{ಕೋಪ} \kn{ಬಂತು}} (\emph{nanage kōpa bantu}) — I got angry: anger \textquotedblleft{}came\textquotedblright{} to me
\end{itemize}

For what happens time after time, \textbf{\kn{ಆಗುತ್ತದೆ}}:

\begin{quote}
\textbf{\kn{ರಾತ್ರಿ} \kn{ನನಗೆ} \kn{ಭಯವಾಗುತ್ತದೆ}.}
\end{quote}

\begin{quote}
\emph{rātri nanage bhayavāguttade.} — \textquotedblleft{}I get scared at night.\textquotedblright{}
\end{quote}

\subsection*{You'll want to know: a feeling you are in}
\begin{quote}
\textbf{\kn{ಅವಳು} \kn{ಈಗ} \kn{ಖುಷಿಯಾಗಿದ್ದಾಳೆ}.}
\end{quote}

\begin{quote}
\emph{avaḷu īga khuṣiyāgiddāḷe.} — \textquotedblleft{}She is happy now.\textquotedblright{}
\end{quote}

\textbf{\kn{ಖುಷಿಯಾಗಿದ್ದಾಳೆ}} is \textbf{\kn{ಖುಷಿ}} + \textbf{\kn{ಆಗಿ}}, \textquotedblleft{}having become glad\textquotedblright{}, + \textbf{\kn{ಇದ್ದಾಳೆ}} (\emph{iddāḷe}, \textquotedblleft{}she is\textquotedblright{}). For \textquotedblleft{}I\textquotedblright{}: \textbf{\kn{ನಾನು} \kn{ಖುಷಿಯಾಗಿದ್ದೇನೆ}} (\emph{nānu khuṣiyāgiddēne}), \textquotedblleft{}I am happy\textquotedblright{}, with the \textbf{\kn{ಇದ್ದೇನೆ}} of \textbf{\kn{ನಾನು} \kn{ಚೆನ್ನಾಗಿದ್ದೇನೆ}}.

Here the person is the subject, in \textbf{\kn{ನಾನು}} or \textbf{\kn{ಅವಳು}}. The \textbf{\kn{ನನಗೆ}} frame says what came over you; the \textbf{-\kn{ಆಗಿದ್ದೇನೆ}} frame says the state you are in. And when you cannot name the feeling, \textbf{\kn{ನನಗೆ} \kn{ಹಾಗೆ} \kn{ಅನಿಸಿತು}} (\emph{nanage hāge anisitu}, \textquotedblleft{}that's how I felt\textquotedblright{}).

\subsection*{Your turn}
Say these aloud:

\begin{itemize}
\raggedright
  \item \textquotedblleft{}Yesterday I was very happy\textquotedblright{} — \emph{ninne nanage tumbā santōṣavāyitu}
  \item \textquotedblleft{}I got angry\textquotedblright{} — \emph{nanage kōpa bantu}
  \item \textquotedblleft{}I get scared at night\textquotedblright{} — \emph{rātri nanage bhayavāguttade}
  \item \textquotedblleft{}I am happy now\textquotedblright{} — \emph{nānu īga khuṣiyāgiddēne}
\end{itemize}

\begin{tcolorbox}[breakable,colback=teal!4,colframe=teal!35!black,title={Before you move on}]
\textquotedblleft{}I was happy\textquotedblright{}? (\textbf{\kn{ನನಗೆ} \kn{ಸಂತೋಷವಾಯಿತು}}.) \textquotedblleft{}I got scared\textquotedblright{}? (\textbf{\kn{ನನಗೆ} \kn{ಭಯವಾಯಿತು}}.) \textquotedblleft{}I got angry\textquotedblright{}? (\textbf{\kn{ನನಗೆ} \kn{ಕೋಪ} \kn{ಬಂತು}}.) \textquotedblleft{}She is happy\textquotedblright{}? (\textbf{\kn{ಅವಳು} \kn{ಖುಷಿಯಾಗಿದ್ದಾಳೆ}}.)
\end{tcolorbox}
\section[\texorpdfstring{\kn{ಗುರಿ} (guri)}{ಗುರಿ (guri)}]{\kn{ಗುರಿ} (guri) — a goal, an aim}

\label{lesson:KA-C315-guri}



\begin{quote}
Before the new one: say \textquotedblleft{}a wish\textquotedblright{} and \textquotedblleft{}a dream\textquotedblright{}. (\emph{āse}, \emph{kanasu}.) Then \textquotedblleft{}She is happy now.\textquotedblright{} (\emph{avaḷu īga khuṣiyāgiddāḷe}.)
\end{quote}

\subsection*{You'll want to know: \kn{ಗುರಿ}}
\textbf{\kn{ಗುರಿ}} — \emph{guri} — \textquotedblleft{}a goal, an aim\textquotedblright{}.

\begin{quote}
\textbf{\kn{ಈ} \kn{ವರ್ಷ} \kn{ನನಗೆ} \kn{ಒಂದು} \kn{ಗುರಿ} \kn{ಇದೆ}: \kn{ಪ್ರತಿದಿನ} \kn{ಕನ್ನಡ} \kn{ಪತ್ರಿಕೆ} \kn{ಓದಬೇಕು}.}
\end{quote}

\begin{quote}
\emph{ī varṣa nanage ondu guri ide: pratidina kannaḍa patrike ōdabēku.} — \textquotedblleft{}This year I have a goal: I must read a Kannada newspaper every day.\textquotedblright{}
\end{quote}

\textbf{\kn{ವರ್ಷ}} (\emph{varṣa}, \textquotedblleft{}a year\textquotedblright{}). \textbf{\kn{ನನಗೆ} … \kn{ಇದೆ}} is \textquotedblleft{}I have\textquotedblright{}, the way \textbf{\kn{ನನಗೆ} \kn{ಜ್ವರ} \kn{ಇದೆ}} is \textquotedblleft{}I have a fever\textquotedblright{}.

\textbf{\kn{ಗುರಿ}} is also the mark an archer aims at, and the aim itself:

\begin{quote}
\textbf{\kn{ಅವನ} \kn{ಗುರಿ} \kn{ತಪ್ಪಿತು}.}
\end{quote}

\begin{quote}
\emph{avana guri tappitu.} — \textquotedblleft{}His aim missed.\textquotedblright{}
\end{quote}

\textbf{\kn{ಅವನ}} (\emph{avana}, \textquotedblleft{}his\textquotedblright{}); \textbf{\kn{ತಪ್ಪಿತು}}, \textquotedblleft{}it slipped, it went astray\textquotedblright{}, the verb of \textbf{\kn{ಬಸ್} \kn{ತಪ್ಪಿಹೋಯಿತು}}.

\begin{cousinweb}
\textbf{\kn{ಗುರಿ}} (\emph{guri}) is native Dravidian. DEDR 1847 lists Kannada \emph{kuṟi} and \emph{guṟi}, \textquotedblleft{}a mark, sign, aim, object of aim\textquotedblright{}, in an older spelling, with a letter for \emph{ṟ} that modern Kannada no longer writes; beside it are Tamil \emph{kuṟi}, \textquotedblleft{}mark, aim, goal\textquotedblright{}, and Telugu \emph{guṟi}, \textquotedblleft{}an aim\textquotedblright{}. A mark you aim at became the goal you aim for.
\end{cousinweb}

\subsection*{Your turn}
Say these aloud:

\begin{itemize}
\raggedright
  \item \emph{guri}, slowly, then at speed
  \item \textquotedblleft{}This year I have a goal\textquotedblright{} — \emph{ī varṣa nanage ondu guri ide}
  \item \textquotedblleft{}His aim missed\textquotedblright{} — \emph{avana guri tappitu}
\end{itemize}

\begin{tcolorbox}[breakable,colback=teal!4,colframe=teal!35!black,title={Before you move on}]
What does \textbf{\kn{ಗುರಿ}} mean? (\textbf{A goal, an aim}; also the mark you aim at.) \textquotedblleft{}I have a goal\textquotedblright{}? (\textbf{\kn{ನನಗೆ} \kn{ಒಂದು} \kn{ಗುರಿ} \kn{ಇದೆ}}.) Is it a native word or a loan? (\textbf{Native Dravidian}: Tamil and Telugu have it too.)
\end{tcolorbox}
\section[\texorpdfstring{\kn{ಹೋಗಬೇಕು} \kn{ಅಂತ} \kn{ಆಸೆ} \kn{ಇದೆ} (hōgabēku anta …)}{ಹೋಗಬೇಕು ಅಂತ ಆಸೆ ಇದೆ (hōgabēku anta …)}]{\kn{ಹೋಗಬೇಕು} \kn{ಅಂತ} \kn{ಆಸೆ} \kn{ಇದೆ} (hōgabēku anta āse ide) — wishes and hopes}

\label{lesson:KA-C315-hogabeku-anta-ase}



\begin{quote}
Say \textquotedblleft{}a goal\textquotedblright{} (\emph{guri}), \textquotedblleft{}I must go\textquotedblright{} (\emph{nānu hōgabēku}) and the word that closes a quotation (\emph{anta}). Today all three go into one sentence.
\end{quote}

\subsection*{You'll want to know: -\kn{ಬೇಕು} \kn{ಅಂತ} \kn{ಆಸೆ} \kn{ಇದೆ}}
\begin{quote}
\textbf{\kn{ನನಗೆ} \kn{ಒಂದು} \kn{ದಿನ} \kn{ಮೈಸೂರಿಗೆ} \kn{ಹೋಗಬೇಕು} \kn{ಅಂತ} \kn{ಆಸೆ} \kn{ಇದೆ}.}
\end{quote}

\begin{quote}
\emph{nanage ondu dina maisūrige hōgabēku anta āse ide.} — \textquotedblleft{}I would love to go to Mysuru one day.\textquotedblright{}
\end{quote}

Taken apart: \textquotedblleft{}to me there is a wish (that says) 'I must go to Mysuru one day'\textquotedblright{}. The wish is a little quotation: the \textbf{-\kn{ಬೇಕು}} of \textquotedblleft{}must\textquotedblright{}, closed by \textbf{\kn{ಅಂತ}}, then \textbf{\kn{ಆಸೆ} \kn{ಇದೆ}}, \textquotedblleft{}there is a wish\textquotedblright{}. In writing you will meet \textbf{\kn{ಬೇಕೆಂದು}} in place of \textbf{\kn{ಬೇಕು} \kn{ಅಂತ}}: \textbf{\kn{ಎಂದು}} is the written twin of \textbf{\kn{ಅಂತ}}.

Swap \textbf{\kn{ಆಸೆ}} for \textbf{\kn{ಕನಸು}}, \textquotedblleft{}a dream\textquotedblright{}, and it becomes a dream: \textbf{… \kn{ಅಂತ} \kn{ಕನಸು} \kn{ಇದೆ}}. For a wish you had then, \textbf{\kn{ಇತ್ತು}}: \textbf{\kn{ನನಗೆ} \kn{ವಿಮಾನದಲ್ಲಿ} \kn{ಹೋಗಬೇಕು} \kn{ಅಂತ} \kn{ಆಸೆ} \kn{ಇತ್ತು}} (\emph{nanage vimānadalli hōgabēku anta āse ittu}, \textquotedblleft{}I used to long to go by plane\textquotedblright{}).

\subsection*{You'll want to know: what you want to become}
\begin{quote}
\textbf{\kn{ನನ್ನ} \kn{ತಂಗಿಗೆ} \kn{ವೈದ್ಯೆ} \kn{ಆಗಬೇಕು} \kn{ಅಂತ} \kn{ಕನಸು} \kn{ಇದೆ}.}
\end{quote}

\begin{quote}
\emph{nanna taṅgige vaidye āgabēku anta kanasu ide.} — \textquotedblleft{}My younger sister dreams of becoming a doctor.\textquotedblright{}
\end{quote}

\textbf{\kn{ತಂಗಿಗೆ}} (\emph{taṅgige}, \textquotedblleft{}to my younger sister\textquotedblright{}); \textbf{\kn{ವೈದ್ಯೆ}} (\emph{vaidye}, \textquotedblleft{}a woman doctor\textquotedblright{}), \textbf{\kn{ವೈದ್ಯ}} with \textbf{-\kn{ಎ}}. \textbf{\kn{ಆಗಬೇಕು}} is \textbf{\kn{ಆಗು}}, \textquotedblleft{}to become\textquotedblright{}, with \textbf{-\kn{ಬೇಕು}}: \textquotedblleft{}must become\textquotedblright{}. The job you want to be needs no ending:

\begin{quote}
\textbf{\kn{ನನಗೆ} \kn{ಶಿಕ್ಷಕ} \kn{ಆಗಬೇಕು} \kn{ಅಂತ} \kn{ಆಸೆ} \kn{ಇದೆ}.}
\end{quote}

\begin{quote}
\emph{nanage śikṣaka āgabēku anta āse ide.} — \textquotedblleft{}I would like to become a teacher.\textquotedblright{}
\end{quote}

\subsection*{Your turn}
Say these aloud:

\begin{itemize}
\raggedright
  \item \textquotedblleft{}I would love to go to Mysuru one day\textquotedblright{} — \emph{nanage ondu dina maisūrige hōgabēku anta āse ide}
  \item \textquotedblleft{}My younger sister dreams of becoming a doctor\textquotedblright{} — \emph{nanna taṅgige vaidye āgabēku anta kanasu ide}
  \item \textquotedblleft{}I would like to become a teacher\textquotedblright{} — \emph{nanage śikṣaka āgabēku anta āse ide}
  \item one place you would love to go, and one thing you would like to become
\end{itemize}

\begin{tcolorbox}[breakable,colback=teal!4,colframe=teal!35!black,title={Before you move on}]
\textquotedblleft{}I would love to go\textquotedblright{}? (\textbf{\kn{ನನಗೆ} \kn{ಹೋಗಬೇಕು} \kn{ಅಂತ} \kn{ಆಸೆ} \kn{ಇದೆ}}.) Which word closes the wish? (\textbf{\kn{ಅಂತ}}.) \textquotedblleft{}To become\textquotedblright{}? (\textbf{\kn{ಆಗಬೇಕು}}.) The written twin of \textbf{\kn{ಬೇಕು} \kn{ಅಂತ}}? (\textbf{\kn{ಬೇಕೆಂದು}}.)
\end{tcolorbox}
\section[(review)]{Review — feelings and hopes}

\label{lesson:KA-R315-feelings-and-hopes}



\begin{quote}
Read \textbf{\kn{ಔಷಧ}} and \textbf{\kn{ಸ್ವಚ್ಛ}} aloud. (\emph{auṣadha}, \emph{svaccha}.) Which letter opens the first, and which sits small inside the second? (\textbf{\kn{ಔ}}; \textbf{\kn{ಛ}}.)
\end{quote}

\subsection*{Your turn}
\begin{itemize}
\raggedright
  \item \emph{Say it:} I was very glad to see you.
  \item \emph{Say it:} When the bus got away, I got angry. Then I felt low.
  \item \emph{Say it:} I am happy now.
  \item \emph{Say it:} This year I have a goal: I would love to go to Mysuru.
  \item \emph{Say it:} My younger sister dreams of becoming a teacher.
  \item \emph{Write it:} \textbf{\kn{ಖುಷಿ}} and \textbf{\kn{ಗುರಿ}}
\end{itemize}

\begin{tcolorbox}[breakable,colback=teal!4,colframe=teal!35!black,title={Before you move on}]
\textquotedblleft{}I was glad\textquotedblright{}? (\textbf{\kn{ನನಗೆ} \kn{ಖುಷಿ} \kn{ಆಯಿತು}}.) \textquotedblleft{}I got scared\textquotedblright{}? (\textbf{\kn{ನನಗೆ} \kn{ಭಯವಾಯಿತು}}.) \textquotedblleft{}I am happy\textquotedblright{}? (\textbf{\kn{ನಾನು} \kn{ಖುಷಿಯಾಗಿದ್ದೇನೆ}}.) \textquotedblleft{}A goal\textquotedblright{}? (\textbf{\kn{ಗುರಿ}}.) \textquotedblleft{}I would love to go\textquotedblright{}? (\textbf{\kn{ಹೋಗಬೇಕು} \kn{ಅಂತ} \kn{ಆಸೆ} \kn{ಇದೆ}}.)
\end{tcolorbox}
